% Generated draft: requires mathematical and source review.
\documentclass[11pt]{article}
\usepackage[margin=25mm]{geometry}
\usepackage{fontspec}
\setmainfont{Latin Modern Roman}
\usepackage{amsmath,amssymb,mathtools,unicode-math}
\setmathfont{Latin Modern Math}
\usepackage{enumitem,needspace,xurl,hyperref}
\setlength{\parindent}{0pt}
\setlength{\parskip}{5pt}
\setlength{\emergencystretch}{3em}
\newcommand{\N}{\mathbb{N}}
\newcommand{\Z}{\mathbb{Z}}
\newcommand{\Q}{\mathbb{Q}}
\newcommand{\R}{\mathbb{R}}
\newcommand{\C}{\mathbb{C}}
\newcommand{\F}{\mathbb{F}}
\newcommand{\A}{\mathbb{A}}
\newcommand{\PP}{\mathbb{P}}
\newcommand{\E}{\mathbb{E}}
\newcommand{\eps}{\varepsilon}
\newcommand{\dd}{\,\mathrm{d}}
\begin{document}
\section*{UnsolvedMath ID: 30005695}
\textbf{Source problem code:} catalogue entry
\subsection{Definitions and mathematical statement}
Let \((W,S)\) be a Coxeter system, so that
\[
W=\langle S\mid (st)^{m_{st}}=1\text{ whenever }m_{st}<\infty\rangle,
\]
where \(m_{ss}=1\) and \(m_{st}=m_{ts}\in\{2,3,\ldots\}\cup\{\infty\}\) for distinct \(s,t\in S\). An entry \(m_{st}=\infty\) imposes no relation on \(st\).

Regard \(W\) as a discrete group and write \(\Q W\) for its group algebra. The rational cohomological dimension \(\operatorname{cd}_{\Q}(W)\) is the projective dimension of the trivial left \(\Q W\)-module \(\Q\). The group is of type \(\mathrm{FP}\) over \(\Q\) if \(\Q\) has a finite resolution by finitely generated projective left \(\Q W\)-modules. In the standard discrete-group sense, a rational duality group of dimension \(d\) is a group of type \(\mathrm{FP}\) over \(\Q\) such that
\[
\operatorname{cd}_{\Q}(W)=d,\qquad
H^i(W,\Q W)=0\quad\text{for every }i\geq0\text{ with }i\neq d,
\qquad H^d(W,\Q W)\neq0.
\]
Here \(H^i(W,M)=\operatorname{Ext}^i_{\Q W}(\Q,M)\), with \(\Q W\) carrying its left regular action. No requirement that \(H^d(W,\Q W)\) be one-dimensional is imposed.

\textbf{Classification problem.} Determine all crystallographic Coxeter groups \(W\) that satisfy these rational duality conditions for some integer \(d\geq1\), and determine the corresponding dimension \(d\).

\textbf{Unresolved scope qualification.} The supplied record does not define ``crystallographic'' or specify a rank convention. A lattice-based interpretation uses a faithful real reflection realization preserving a full lattice, meaning a discrete subgroup spanning its finite-dimensional real representation space. The record does not identify the intended realization or establish that this interpretation is its convention. Thus the rational duality condition above is explicit, but the precise class of Coxeter groups to which the classification applies requires confirmation from the source. No finite-rank, irreducibility, affine, or hyperbolic restriction is added here.
\subsection{Short English statement}
Classify the crystallographic Coxeter groups whose rational group-ring cohomology is nonzero in exactly one degree \(d\geq1\), subject to the finite projective resolution condition stated above. The intended meaning and rank scope of ``crystallographic'' require source confirmation. This is a statement-only draft requiring human mathematical review.
\subsection{Sources}
\begin{itemize}
\item Supplied local UnsolvedMath registry record, permanent ID \(30005695\), source problem code \texttt{catalogue entry}, title ``Rational Duality among Crystallographic Coxeter Groups.'' No individual proposer is supplied. The registry gives proposed year \(2023\).
\item Supplied report citation: \emph{Mini-Workshop: Homological aspects of TDLC-groups} (2024), Oberwolfach Reports, pp.~3103--3138, DOI \texttt{10.4171/owr/2023/55}. \url{https://doi.org/10.4171/owr/2023/55}. The raw excerpt contains the page marker \(3113\).
\item Supplied literature source, ``Oberwolfach report'': \url{https://ems.press/content/serial-article-files/48166}.
\item Supplied literature source, ``Castellano--Weigel'': \url{https://doi.org/10.1016/j.jalgebra.2016.01.008}.
\item Supplied literature source, ``Davis--Kahle'': \url{https://doi.org/10.1112/jtopol/jtt047}.
\item Literature assessment checked at \texttt{2026-08-22}: this is supplied dated curation, not independently verified. It labels the problem open and states that affine and crystallographic hyperbolic groups give examples, while the report records no general classification of crystallographic Coxeter groups as rational duality groups. The supplied background explicitly calls this dated literature triage and requests independent verification of time-sensitive status and cited claims; the example claim is not expanded here into a universal assertion about either class.
\item The registry labels the statement \texttt{corrected\_verified} and says that it was checked against the official report and preserves its hypotheses, notation, and target. These are supplied source-verification assertions, not verification performed for this draft. The record's open status is likewise a supplied label.
\end{itemize}
\end{document}
